% =======================================================
% 2. QUALITA DI PROCESSO
% =======================================================
\section{Qualità di Processo}
Il gruppo adotta lo standard \textbf{ISO/IEC 12207} per la gestione dei processi. La qualità è perseguita attraverso il monitoraggio costante di indicatori quantitativi calcolati al termine di ogni Sprint\textsubscript{G}.

\subsection{Processi primari}
\subsubsection{Fornitura}
Questo processo comprende le attività volte a determinare le procedure per la pianificazione e il controllo del progetto. Il Responsabile\textsubscript{G} monitora l'avanzamento dei lavori e il consumo del budget\textsubscript{G} per garantire trasparenza alla proponente Zucchetti.

\subsubsection{Sviluppo}
Include le attività di analisi, progettazione e codifica. Durante la precedente fase di RTB, il focus è sulla stabilità dei requisiti estratti dal capitolato C6, al fine di minimizzare revisioni costose nelle fasi successive.

\subsection{Processi di supporto}
\subsubsection{Documentazione}
Assicura che ogni documento prodotto sia leggibile e conforme agli standard fissati nelle Norme di Progetto. La qualità è misurata tramite l'indice di leggibilità e l'assenza di errori formali.

\subsubsection{Gestione configurazione}
Garantisce l'integrità dei prodotti di lavoro tramite il versionamento su repository\textsubscript{G} Git. Ogni modifica deve essere tracciabile e associata a una specifica issue\textsubscript{G}.

\subsubsection{Verifica}
Accerta che ogni attività non contenga errori tramite un processo di controllo incrociato tra i membri del gruppo. Per i documenti, la verifica manuale si focalizza su sintassi, semantica e aderenza alle norme.
\vspace{0.5cm}
\subsubsection*{Checklist di Ispezione Documentale}
Per garantire uniformità e oggettività durante le attività di verifica manuale, i Verificatori si avvalgono della seguente lista di controllo strutturata. Questa tabella raggruppa le verifiche fondamentali da effettuare prima di approvare una Pull Request relativa a un documento:

\begin{table}[H]
    \centering
    \renewcommand{\arraystretch}{1.3}
    \begin{tabularx}{\textwidth}{|l|X|}
        \hline
        \rowcolor{primarycolor!10}
        \textbf{Elemento} & \textbf{Descrizione del Controllo} \\
        \hline
        \textbf{Dati di Copertina} & Verifica della corretta corrispondenza tra versione, data, autori e verificatori nel Registro delle Modifiche e nel frontespizio. \\
        \hline
        \textbf{Terminologia (Glossario)} & Controllo che i termini tecnici, alla loro prima occorrenza nel testo, siano rigorosamente contrassegnati dal pedice\textsubscript{G}. \\
        \hline
        \textbf{Didascalie e Riferimenti} & Accertamento che ogni figura e tabella inserita possieda una \textit{caption} esplicativa e sia correttamente richiamata nel testo. \\
        \hline
        \textbf{Struttura} & Assenza di sezioni vuote prive di contenuto o di gerarchie dei titoli errate. \\
        \hline
        \textbf{Leggibilità (Gulpease)} & Rilevazione di frasi eccessivamente lunghe o complesse per garantire il rispetto della soglia ottimale, con l'ausilio di script automatici. \\
        \hline
        \textbf{Tracciamento} (AdR / PdQ) & Verifica della bidirezionalità dei tracciamenti (es. Caso d'Uso $\leftrightarrow$ Requisito $\leftrightarrow$ Test di Sistema). \\
        \hline
        \textbf{Ortografia e Formattazione} & Controllo sintattico e semantico per eliminare refusi, doppi spazi e disallineamenti tipografici rispetto alle Norme di Progetto. \\
        \hline
    \end{tabularx}
    \caption{Checklist per l'ispezione formale della documentazione}
\end{table}

\subsubsection{Validazione}
Fornisce la conferma che il prodotto finale sia conforme alle attese dell'utente e della proponente. In questa fase si concentra sulla verifica che i requisiti individuati rispondano agli obiettivi del "Second Brain".

\subsubsection{Accertamento della qualità}
Monitora l'efficacia dei processi adottati e il raggiungimento degli obiettivi qualitativi definiti nel presente documento. Tale processo si basa sulla raccolta periodica delle metriche, sull'analisi degli scostamenti e sull'individuazione di eventuali azioni correttive.

\subsubsection{Gestione dei cambiamenti}
Disciplina la valutazione, l'approvazione e il tracciamento delle modifiche rilevanti a requisiti, documenti, tecnologie e processi. Ogni cambiamento significativo viene gestito tramite issue\textsubscript{G}, Pull Request\textsubscript{G} e verifica da parte di un membro diverso dall'autore.

\subsubsection{Gestione delle non conformità}
Gestisce le deviazioni rispetto alle Norme di Progetto, al Piano di Qualifica o agli accordi formalizzati nei verbali. Le non conformità vengono classificate, tracciate e risolte tramite attività correttive dedicate.

\subsection{Processi organizzativi}
\subsubsection{Gestione processo}
Si occupa del coordinamento del team durante i cicli di Sprint. Mira a bilanciare il carico di lavoro tra i componenti del gruppo Synergo Unit per evitare colli di bottiglia\textsubscript{G}.

\subsubsection{Miglioramento continuo}
Attraverso le retrospettive di fine Sprint, il gruppo analizza le inefficienze e definisce azioni correttive immediate per ottimizzare le metodologie di lavoro.

\subsubsection{Comunicazione}
Regola i canali e le modalità di comunicazione interna ed esterna, garantendo che decisioni, chiarimenti e impegni siano tracciati tramite verbali, issue o comunicazioni ufficiali.

\subsubsection{Infrastruttura}
Definisce e mantiene gli strumenti utilizzati dal gruppo per documentazione, pianificazione, sviluppo del Proof of Concept\textsubscript{G}, pubblicazione e tracciamento delle attività.

\subsubsection{Formazione}
Supporta l'acquisizione progressiva delle competenze necessarie all'utilizzo degli strumenti, delle tecnologie e dei processi adottati. In fase RTB assume particolare rilevanza per mitigare i rischi legati all'inesperienza tecnica.

\subsection{Metriche per processi}
Le metriche di processo (MPC) vengono calcolate a ogni fine Sprint e riportate nel Cruscotto delle Metriche.
Le metriche di processo sono inoltre utilizzate come strumenti di monitoraggio dei rischi individuati nel 
\link{https://synergo-unit-unipd.github.io/Second-Brain/docs/PB/doc_gestionali/doc_esterni/piano_di_progetto/piano_di_progetto.pdf}{Piano di Progetto}. 
In particolare, SPI e CPI supportano il controllo dei rischi legati a scadenze e budget; RSI consente di monitorare la stabilità dei requisiti; 
Time Efficiency e Task Completion Rate aiutano a rilevare criticità organizzative; PR Resolution Time e Review Coverage permettono di controllare 
l'efficacia del processo di verifica; Quality Metrics Compliance Rate fornisce una visione sintetica del livello complessivo di qualità raggiunto.

\vspace{0.5cm}
\subsubsection*{Definizione matematica delle metriche principali}
Di seguito sono riportate le formule matematiche utilizzate per il calcolo degli indici di performance economica, temporale e di stabilità, ricavate dal foglio di tracciamento Excel:

\begin{itemize}

    \item \textbf{Schedule Performance Index (SPI)}: Misura l'efficienza temporale del progetto. Un valore $\geq 1$ indica che il progetto è in linea con la pianificazione o in anticipo.

    \[
    SPI = \frac{EV}{PV}
    \]

    dove:
    \begin{itemize}
        \item $EV$ (\textit{Earned Value}) rappresenta il valore del lavoro effettivamente completato;
        \item $PV$ (\textit{Planned Value}) rappresenta il valore del lavoro pianificato.
    \end{itemize}

    \item \textbf{Cost Performance Index (CPI)}: Misura l'efficienza economica del progetto. Un valore $\geq 1$ indica che il costo sostenuto è inferiore o uguale al budget previsto per il lavoro svolto.

    \[
    CPI = \frac{EV}{AC}
    \]

    dove:
    \begin{itemize}
        \item $EV$ (\textit{Earned Value}) rappresenta il valore del lavoro effettivamente completato;
        \item $AC$ (\textit{Actual Cost}) rappresenta il costo effettivamente sostenuto.
    \end{itemize}

    \item \textbf{Requirements Stability Index (RSI)}: Misura la stabilità dei requisiti consolidati.

    \[
    RSI = \left( 1 - \frac{N_{modifiche}}{N_{totali}} \right) \times 100
    \]

    dove:
    \begin{itemize}
        \item $N_{modifiche}$ rappresenta il numero di requisiti aggiunti, modificati o rimossi nel periodo considerato;
        \item $N_{totali}$ rappresenta il numero totale di requisiti presenti all'inizio del periodo considerato.
    \end{itemize}

    \item \textbf{Quality Metrics Compliance Rate (QMCR)}: Misura la percentuale di metriche che rispettano i valori soglia definiti.

    \[
    QMCR = \frac{N_{conformi}}{N_{misurate}} \times 100
    \]

    dove:
    \begin{itemize}
        \item $N_{conformi}$ rappresenta il numero di metriche che rispettano i valori soglia previsti;
        \item $N_{misurate}$ rappresenta il numero totale di metriche effettivamente misurate nel periodo considerato.
    \end{itemize}

\end{itemize}
\vspace{0.5cm}

\begin{table}[H]
    \centering
    \renewcommand{\arraystretch}{1.3}
    \begin{tabularx}{\textwidth}{|l|X|c|c|}
        \hline
        \rowcolor{primarycolor!10}
        \textbf{Codice} & \textbf{Nome Metrica} & \textbf{Val. Accettabile} & \textbf{Val. Ottimale} \\
        \hline
        MPC01 & \textbf{Schedule Performance Index\textsubscript{G} (SPI)}: efficienza nel rispettare i tempi pianificati. & $\ge 0.9$ & $1.0$ \\
        \hline
        MPC02 & \textbf{Cost Performance Index\textsubscript{G} (CPI)}: efficienza nell'uso del budget allocato. & $\ge 0.9$ & $1.0$ \\
        \hline
        MPC03 & \textbf{Requirements Stability Index\textsubscript{G} (RSI)}: variazione dei requisiti nell'Analisi. & $\ge 80\%$ & $100\%$ \\
        \hline
        MPC04 & \textbf{Time Efficiency}: percentuale di ore produttive sul totale ore spese. & $\ge 75\%$ & $\ge 90\%$ \\
        \hline
        MPC05 & \textbf{Task\textsubscript{G} Completion Rate}: percentuale di issue chiuse rispetto a quelle pianificate. & $\ge 80\%$ & $100\%$ \\
        \hline
        MPC06 & \textbf{PR Resolution Time}: tempo medio per revisione\textsubscript{G} e merge\textsubscript{G} di una Pull Request\textsubscript{G}. & $\le 3 \text{ giorni}$ & $\le 1 \text{ giorno}$ \\
        \hline
        MPC07 & \textbf{Review Coverage}: percentuale di Pull Request approvate da un revisore diverso dall'autore. & $100\%$ & $100\%$ \\
        \hline
        MPC08 & \textbf{Quality Metrics Compliance Rate}: percentuale di metriche i cui valori sono almeno accettabili. & $\ge 80\%$ & $\ge 95\%$ \\
        \hline
    \end{tabularx}
    \caption{Metriche di Qualità di Processo}
\end{table}

\vspace{0.5cm}
\subsubsection*{Motivazione delle soglie metriche}
\begin{itemize}
    \item \textbf{MPC01 (SPI) e MPC02 (CPI)}: Un valore accettabile $\ge 0.9$ tollera un fisiologico margine di errore del 10\% nelle stime, tipico di un team con esperienza in maturazione. Il valore ottimale $1.0$ rappresenta la perfetta aderenza al piano e al budget.
    \item \textbf{MPC03 (RSI)}: La soglia accettabile dell'80\% assorbe i naturali cambiamenti emersi durante i colloqui di chiarimento con la proponente in fase RTB. Un valore ottimale del 100\% indica un'analisi iniziale perfetta e assenza di ambiguità.
    \item \textbf{MPC04 (Time Efficiency)}: Si accetta un $\ge 75\%$ poiché una parte del tempo deve necessariamente essere investita in auto-formazione ("ore di palestra") e coordinamento. Un valore $\ge 90\%$ indica un team pienamente a regime e con basso overhead organizzativo.
    \item \textbf{MPC05 (Task Completion Rate)}: La soglia dell'80\% tollera che alcune task possano slittare allo sprint successivo per impedimenti non prevedibili. Il 100\% ottimale indica una pianificazione dello sprint impeccabile.
    \item \textbf{MPC06 (PR Resolution Time)}: Un tempo accettabile di $\le 3$ giorni previene la stagnazione del codice e riduce il rischio di complessi conflitti di merge. L'ottimo $\le 1$ giorno favorisce una vera Continuous Integration.
    \item \textbf{MPC07 (Review Coverage)}: Il 100\% è l'unico valore ammesso, sia come accettabile che come ottimale, poiché ogni Pull Request deve essere approvata da un revisore diverso dall'autore prima di poter essere integrata nel ramo principale.
    \item \textbf{MPC08 (Quality Metrics Compliance)}: Un'accettazione all'80\% riconosce che singole metriche possano temporaneamente fallire durante uno sprint, pur mantenendo la qualità generale sotto controllo.
\end{itemize}